\section{Monica：从浏览器插件到 AI 应用平台}

前面讲第一段创业方法，本章进入第二段创业。2022 年底，肖弘团队看到生成式 AI 的潜力，开始做 Monica。最初选择浏览器插件，而不是独立 chatbot app，是因为插件更容易嵌入用户已有工作流；后来收购 ChatGPT for Google，又给 Monica 提供了冷启动流量和用户反馈。

\begin{figure}[H]
\centering
\includegraphics[width=0.96\textwidth]{monica-product-evolution.png}
\caption{Monica 产品演进：从浏览器插件、ChatGPT for Google 到 AI 应用平台。自制概念图，依据 01:05:55--01:48:38 对谈内容整理。}
\end{figure}

\begin{knowledgebox}{读图：插件不是终局，但它是进入工作流的入口}
从 Plugin 到 ChatGPT for Google、Monica.im、用户反馈和 Agent，图里强调的是冷启动和工作流嵌入。插件让 AI 出现在用户已经工作的地方，而不是要求用户主动打开另一个聊天窗口。
\end{knowledgebox}

\subsection{为什么不直接做 Chatbot App}

本节解释产品形态选择。独立 chatbot app 要回答“用户为什么打开你而不是 ChatGPT 原厂”的问题；插件可以在搜索、网页阅读、写作、翻译、总结等场景中出现，降低启动阻力。肖弘也承认，ChatGPT 官方 App 发布前存在一个 app 窗口期，一些海外套壳产品抓到了，但 Monica 的选择是以插件切入工作流。

\begin{knowledgebox}{术语消化：冷启动、工作流、容器}
\noindent\begin{tabularx}{\textwidth}{p{0.18\textwidth}p{0.36\textwidth}X}
\toprule
术语 & 含义 & 本期中的作用 \\
\midrule
冷启动 & 产品早期缺用户和反馈时的启动问题 & ChatGPT for Google 给 Monica 带来用户和反馈。 \\
工作流 & 用户完成任务的连续步骤和工具环境 & 插件能嵌入已有工作流。 \\
容器 & 承接模型能力的产品形态 & 好容器把模型能力变成可用任务。 \\
用户理解 & 对不同地区、文化和使用方式的深入理解 & 出海产品不能只翻译界面。 \\
\bottomrule
\end{tabularx}
\end{knowledgebox}

\subsection{为什么重新做 Monica，而不是改名 ChatGPT for Google}

本节保留一个关键产品判断。ChatGPT for Google 有自然增长和品牌词，但它更像功能名称，不像长期品牌。直接改名可能伤害增长；继续只维护旧产品，团队也学不到海外增长、投放、KOL 和用户理解。于是团队重新做 Monica，同时保留旧产品作为安全垫。这样既有冷启动资产，又让团队获得真实增长训练。

\begin{importantbox}{课堂提示：买来的增长不能替代团队学习}
并购一个有流量的产品，可以缩短冷启动时间；但如果团队只躺在自然增长上，就学不到增长和用户理解。重新做 Monica 的意义，是把团队放回需要自己获取反馈的环境。
\end{importantbox}

\subsection{出海：勇气先于技巧}

本节看海外市场选择。肖弘认为，国内是统一大市场，赢了很大但非常卷；海外更像许多小山头，语言、货币、时区、审美和文化都不同，给创业公司留下更多局部空间。Monica 选择先做海外，需要勇气：团队没有长期海外生活经历，英语也不是核心优势，但窗口期比完美准备更重要。

\begin{warningbox}{出海不是翻译界面}
海外产品需要考虑多语言、多货币、时区、支付、合规、审美和用户习惯。AI 可以帮助翻译，但不能自动解决用户理解。真正的出海，是组织持续学习不同市场。
\end{warningbox}

\subsection{本章小结}

Monica 的路径说明，AI 应用创业不是简单套模型，而是选择合适入口、冷启动方式、品牌名称、市场范围和反馈机制。它把第一段创业的商业化、增长和平台预判经验，迁移到了 AI 应用窗口期。

